\documentclass[letterpaper,11pt]{article}
\usepackage[top=0.8in, bottom=0.8in, left=0.9in, right=0.9in]{geometry}
\usepackage[hidelinks]{hyperref}
\usepackage{lmodern}
\usepackage{parskip}
\usepackage{microtype}
\setlength{\parskip}{6pt}
\setlength{\parindent}{0pt}
\begin{document}
\begin{flushleft}
\textbf{\large Ajay Venkatesh} \\
Brooklyn, NY \\
+1 516 585 9013 \\
\href{mailto:av3855@nyu.edu}{av3855@nyu.edu}
\end{flushleft}
\vspace{10pt}
May 21, 2026
\vspace{6pt}

Hiring Manager \\
Lightning AI

\vspace{10pt}

Dear Hiring Manager,

My name is Ajay Venkatesh. I'm finishing my M.S. in Computer Engineering at NYU, with production software experience at PwC in Bengaluru, a summer SDE internship at Amazon, and ongoing on-campus work at NYU's Novel AI Technologies. I'm writing about the Fullstack Engineer role on the Lightning AI platform.

The role's framing, building an end-to-end platform that spans frontend, CLI, APIs, and core backend for AI workloads, sits exactly where I've been gravitating. At Amazon I deployed a \textbf{Retrieval-Augmented Generation (RAG)} pipeline with an \textbf{agentic workflow} that autonomously reasons over operational logs, and implemented a \textbf{Model Context Protocol (MCP)} server for tool-augmented retrieval that materially reduced query time and compute costs. On top of that I provisioned resilient \textbf{Infrastructure-as-Code} on \textbf{AWS CDK} across auto-scaling ECS clusters with multi-stage CI/CD and rollback controls. Owning AI-adjacent features end-to-end across infra and product is the shape of work the posting describes.

On full-stack and frontend: at NYU I shipped a customer-facing analytics platform end-to-end with \textbf{React, Node.js, REST APIs}, deployed via \textbf{Docker} on \textbf{EC2}, with Stripe billing, role-based access, and row-level security for paying clients. At Amazon I built a \textbf{React 18} interface with optimized API integration and client-side caching. At Campus Mesh I shipped \textbf{Node.js microservices} powering real-time academic workflows, plus an LLM-powered AI assistant with prompt engineering and RAG retrieval.

Stack-wise: \textbf{React}, \textbf{Python}, \textbf{Node.js}, \textbf{TypeScript}, \textbf{Docker}, \textbf{Kubernetes}, \textbf{AWS}, plus daily work on \textbf{PyTorch}, \textbf{LangChain}, and \textbf{Hugging Face}. I'm comfortable maintaining high standards for code quality, implementing automation, and reducing time to production while keeping technical debt under control.

What pulls me toward Lightning AI is the leverage. PyTorch Lightning shapes how a huge fraction of ML practitioners actually train and deploy models, and getting the platform layer right compounds across thousands of teams. That's the kind of platform engineering where the work matters far beyond the immediate user.

I'd welcome the chance to discuss further. Thanks for considering my application.

\vspace{12pt}
Sincerely, \\
Ajay Venkatesh

\end{document}